\section{Method}
\label{sec:method}

\subsection{Task, world pairs, and two moments}
\label{sec:task}

\paragraph{Task.}
An episode starts at a viewpoint and heading $s_0$ in a simulated scene.
The instruction is ``Select the marked object closest to the $\langle$anchor$\rangle$.''
Two candidate objects $A$ and $B$ carry numeric marks~\cite{som2023} whenever they are visible and within reach; the anchor $X$ is an object of a different category that is unique in the scene.
The correct target $g(w)$ in world $w$ is the candidate whose center is closer to $X$, and we require the farther candidate to be at least $1.5$ times as far, so that the answer is unambiguous.
At each turn the agent receives the current frame, the history, and a text listing the available actions.
It can take an \emph{observation action}, \ie turn left or right by $45^\circ$ or move to an adjacent viewpoint, each costing $\lambda$, or it can end the episode with \commit$(c)$ for a marked candidate $c$, which yields reward $\mathbb{1}[c = g(w)]$.
Marks are placed only on candidates that can legitimately be selected, so the reward equals task success.
An episode allows at most $D$ observation actions and $D{+}1$ turns.

\paragraph{World pairs and cases.}
A world pair $(w_0, w_1)$ shares the scene, the start $s_0$, the candidates, and the instruction, and differs only in the pose of $X$; all other object poses are copied from $w_0$ (\cref{sec:data}).
Two binary properties define four cases (\cref{tab:cases}).
In \emph{direct} pairs $X$ is visible from $s_0$ in both worlds, so the evidence is in the input from the start; in \Look\ pairs $X$ is out of view at $s_0$ in both worlds, the start image is shared, and one observation action reveals it.
In \Flip\ pairs $g(w_0) \neq g(w_1)$; in \Stay\ pairs $g(w_0) = g(w_1)$.
\Look$\to$\Stay\ pairs prevent ``looking means switching'' from being a valid shortcut.

\begin{table}[t]
  \centering
  \small
  \begin{tabular}{@{}lcc@{}}
    \toprule
    & target unchanged & target changes \\
    \midrule
    anchor visible at $s_0$ & direct \Stay & direct \Flip \\
    anchor needs one look & \Look$\to$\Stay & \Look$\to$\Flip \\
    \bottomrule
  \end{tabular}
  \caption{The four cases. Labels are defined relative to what the agent has observed, not to the world alone.}
  \label{tab:cases}
\end{table}

\paragraph{Two moments.}
Let $h_t(w)$ be the agent's full input at turn $t$ in world $w$.
We use two moments that are easy to conflate.
$\tau$ is the first turn at which $h_t(w_0) \neq h_t(w_1)$; for all $t < \tau$ the inputs are byte-identical, so the two worlds form one \emph{information set}.
$\tauseen(w)$ is the first turn at which the anchor is visible in the agent's own frame in world $w$.
The two can differ: in a \Look$\to$\Flip\ pair the agent may first see the region next to candidate $A$ with no anchor in it.
The inputs now differ, so $\tau$ has passed, but the agent has not seen $X$, and a \commit\ at that point is still a guess.
We use $\tau$ to decide how credit is shared and $\tauseen$ to decide whether a commitment was premature.
Both are computed from the agent's own trajectory and precomputed view tables, so they are defined identically for single-world and paired training.

\subsection{Credit before the evidence}
\label{sec:credit}

\paragraph{Commit-now value.}
For a turn $t$ at which at least one candidate is marked, let
\begin{equation}
\Qc(w, t) = \sum_{c} \pi_\theta(\commit(c) \mid h_t)\, \mathbb{1}[c = g(w)],
\label{eq:qcommit}
\end{equation}
where the probabilities are obtained by scoring the action text of every selectable candidate and renormalizing over them.
This is the expected reward of committing immediately; it requires no additional rollout.
Let $G(w, t)$ be the realized return from turn $t$ in world $w$: the final reward minus $\lambda$ for every later observation action.

\paragraph{Single-world contrast.}
A same-world comparison of looking against committing, in the spirit of~\cite{seekvln2026}, assigns to a pre-evidence turn ($t < \tau$) of an episode in the realized world $w^\ast$ the advantage
\begin{equation}
A^{\mathrm{single}}_t = G(w^\ast, t) - \Qc(w^\ast, t).
\end{equation}
If the agent commits at $t < \tau$ and its guess happens to be right in $w^\ast$, the commitment is credited positively even though the same input would have led to failure in the other world.

\paragraph{Information-set credit.}
With lockstep pairs we observe both outcomes.
While $t < \tau$, the two prompts are byte-identical, so the policy generates one response and it is executed in both worlds; after $\tau$ the two episodes continue independently.
We implement this by deduplicating identical prompts within a pair, which requires no change to the one-environment-one-trajectory structure of the trainer.
For $t < \tau$ we assign to both copies of the turn the same advantage
\begin{equation}
A^{\mathrm{IS}}_t = \sum_{w \in \{w_0, w_1\}} p(w) \big[ G(w, t) - \Qc(w, t) \big],
\label{eq:is}
\end{equation}
where $p(w_1) = \rho$ is the training change rate and $p(w_0) = 1 - \rho$.
After $\tau$, and throughout direct pairs, we use standard multi-turn GAE~\cite{gae2016,vagen2025}.
By construction, two worlds with identical inputs receive identical credit for the same decision.

\paragraph{What pairing changes: a $2{\times}2$ view.}
Under lockstep sampling the shared turn appears in both trajectories and the two gradients add.
Training on paired data with the critic's GAE in each world therefore already merges the two worlds; to first order it differs from \cref{eq:is} only in the baseline used before the evidence.
Information-set credit thus combines two ingredients, which we test separately:
(i) \emph{sampling}: single-world episodes versus lockstep pairs, and
(ii) \emph{pre-evidence baseline}: the learned critic versus $\Qc$.
For the sampling factor, write $X_w = G(w, t) - \Qc(w, t)$ for a continuation sampled in world $w$.
If worlds are drawn from the same prior, single-world and paired estimates have the same expectation.
At an equal budget of two continuations and $\rho = \tfrac12$, the mean of two single-world estimates has variance $\tfrac12\big(\mathbb{E}_w[\mathrm{Var}\,X_w] + \mathrm{Var}_w(\mathbb{E}\,X_w)\big)$, whereas the paired estimate has variance $\tfrac12\,\mathbb{E}_w[\mathrm{Var}\,X_w]$.
Pairing is stratified sampling over the hidden world: it removes the between-world term $\mathrm{Var}_w(\mathbb{E}\,X_w)$, which is large precisely when committing early pays off in one world and fails in the other.
We do not claim that single-world credit is biased.
For $\rho \neq \tfrac12$ one rollout per world is no longer a proportional allocation; we weight each trajectory's loss by $2p(w)$ so that paired data represent the intended world distribution, and treat the variance benefit as an empirical question.

\paragraph{Turns without a valid commit-now value.}
If no candidate is marked at a pre-evidence turn, or the model's output cannot be parsed, $\Qc$ is undefined.
On such turns we fall back to the critic's token-level GAE, averaged over the two worlds for paired methods so that both copies remain identical.
All advantage replacements are done before a single batch-level normalization, and the critic is trained on its usual return targets.

\subsection{Post-evidence consistency}
\label{sec:consistency}
Once the agent has seen the anchor, the correct behavior depends on the case.
We add a consistency term only at the final \commit\ turn of an episode in world $w$, and only if $t \ge \tauseen(w)$; before the agent has seen $X$, forcing a switch would teach it to infer from absence.
Let $s_w(c)$ be the policy's score of \commit$(c)$ at that turn, computed with the same scoring as $\Qc$, and let $\bar w$ denote the other world of the pair.
For \Flip\ pairs we use a twin margin on the log-probabilities,
\begin{equation}
c_w^{\Flip} = \max\!\big(0,\; m - [\, s_w(g(w)) - s_w(g(\bar w)) \,]\big),
\end{equation}
which asks the agent to prefer its own world's target over the other world's by a margin $m$.
For \Stay\ pairs, when both episodes have seen the anchor, we pull the two target distributions together,
\begin{equation}
c_w^{\Stay} = \mathrm{KL}\big(\pi^{\mathrm{sel}}_w \,\|\, \mathrm{sg}[\hat\pi^{\mathrm{sel}}_{\bar w}]\big),
\end{equation}
where $\pi^{\mathrm{sel}}_w$ is the policy's distribution over $\{A, B\}$ at the \commit\ turn, $\hat\pi^{\mathrm{sel}}_{\bar w}$ is the same distribution in the other world computed at sampling time, and $\mathrm{sg}$ stops the gradient.
Each episode carries one direction, so a \Stay\ pair contributes both directions of the KL.
Terms whose candidates are not selectable at that turn are set to zero.
The auxiliary loss is $\mathcal{L}_{\mathrm{cons}} = \frac{1}{N}\sum_i 2p(w_i)\, c_i$ over the $N$ episodes of a PPO mini-batch, and the full objective is the PPO loss plus $\beta\, \mathcal{L}_{\mathrm{cons}}$.
Neither term constrains the route, which may legitimately differ between worlds.

\subsection{Reward and optimization}
\label{sec:training}
All methods use the same reward: $+1$ for selecting the correct target, $0$ otherwise, and $-\lambda$ per observation action.
We do not use VAGEN's format reward, which gives each well-formed action a bonus far larger than $\lambda$ and would itself encourage observation.
We train with the multi-turn PPO~\cite{ppo2017} of VAGEN~\cite{vagen2025} and use a KL penalty to the reference policy as a loss term rather than in the reward, so that replacing pre-evidence advantages leaves the KL term unchanged and identical across methods.
Every method uses the same number of episodes per training step: single-world methods sample $64$ episodes, and paired methods sample $32$ pairs.
